\documentclass[../../proceedings/main.tex]{subfiles}
\begin{document}
\begin{refsection}[\subfix{references.bib}]
\PaperHeading
  {Cross-Platform AI Trend Radar}
  {Alex Rivera and Samira Becker}
  {Faculty of Media, Hochschule D\"usseldorf}
  {trend-radar@example.edu}
  {AI research and product signals appear across papers, code repositories, and public discussion. We present a small cross-platform trend radar that collects these signals, normalizes topic labels, and ranks changes over time. A transparent score combines relative activity, recent growth, and source diversity. In an illustrative evaluation on a hand-labelled sample, the combined ranking surfaces more relevant emerging topics than a volume-only baseline. The prototype demonstrates a reproducible workflow and makes each score traceable to source counts. Its main limitations are small sample size, platform bias, and the instability of short observation windows.}
  {trend detection, information retrieval, AI monitoring, evaluation}

\section{Introduction and Motivation}
New AI topics often emerge in several places at once: research papers describe
methods, code repositories expose implementations, and community discussions
reveal practical interest. A list sorted by raw mentions tends to reward already
popular topics. Our project asks whether a small, explainable score can identify
recent changes while retaining the origin of each signal. The pipeline is
inspired by established work on attention mechanisms \autocite{vaswani2017attention}
and information retrieval evaluation \autocite{manning2008ir}.

\section{Approach}
The prototype ingests weekly counts from three source categories: preprints,
public repositories, and discussion posts. For this example, the counts are a
small synthetic dataset; no live API or personal data is needed for the build.
Topic aliases are mapped to one canonical label, then each platform's counts
are normalized by its weekly total. Figure~\ref{fig:radar} shows the stages.
\begin{figure}[t]
  \centering
  \includegraphics[width=\columnwidth]{\subfix{figures/radar.pdf}}
  \caption{The trend radar combines normalized signals from three platforms.}
  \label{fig:radar}
\end{figure}

For topic $t$, let $g_t$ be normalized week-over-week growth, $a_t$ current
relative activity, and $d_t$ the fraction of source categories with activity.
The illustrative score in Equation~\ref{eq:score} is a weighted average:
\begin{equation}
  S_t = 0.45g_t + 0.35a_t + 0.20d_t.
  \label{eq:score}
\end{equation}
Each term is scaled to $[0,1]$ before aggregation. We retain component values
alongside the final rank to make changes inspectable.

\section{Design Decisions and Alternatives}
The three weights favor emerging activity while keeping a penalty for a topic
seen on only one platform. An unweighted sum was simpler but made persistent,
high-volume topics dominate. A learned ranker might adapt weights, yet our
sample is too small for a credible training and test split. We therefore use
fixed weights and report them explicitly. Deduplication uses canonical topic
labels rather than embeddings, avoiding opaque merges at this stage.

\section{Results}
We manually marked ten candidate topics as relevant or irrelevant to the
exercise. Table~\ref{tab:comparison} reports precision among the top five
ranked topics. The combined score placed four relevant topics in the top five;
the volume-only baseline placed three. These are illustrative observations,
not a general claim about AI trend forecasting.
\begin{table}[t]
  \centering
  \caption{Illustrative ranking on ten hand-labelled topics.}
  \label{tab:comparison}
  \begin{tabular}{lcc}\toprule
    Ranking method & Relevant & Precision@5 \\ \midrule
    Volume-only baseline & 3 & 0.60 \\
    Combined score & 4 & 0.80 \\ \bottomrule
  \end{tabular}
\end{table}
The small difference suggests that growth and source diversity are useful
features in this toy setting. A real deployment would require a larger,
time-separated evaluation.

We also inspected the score components rather than only the final order. One
candidate had high activity on repositories but almost no preprint or
discussion signal. Its diversity term was low, so it ranked below a topic
with moderate activity across all three sources. This behavior is intentional:
the radar seeks broad change, not the busiest individual platform. Another
candidate had a large growth ratio because its earlier count was close to
zero. We capped the normalized growth term at one, preventing a single small
denominator from overwhelming the ranking. These inspection cases help check
whether the score behaves as specified, but they do not establish predictive
validity.

The hand labels were created after the candidate set was assembled. That
procedure can favor topics that are easy to recognize in retrospect. A stronger
study would select evaluation weeks in advance, hide future counts from the
annotators, and compute precision at several cutoffs. It would also record
inter-annotator agreement and compare against a recent-growth baseline. We
report only Precision@5 here because the ten-topic exercise is too small for
stable conclusions about recall or long-term forecasting.

\section{Discussion and Limitations}
Platform coverage determines what the radar can see. Missing APIs, changing
moderation practices, and language imbalance can shift counts independently of
actual research interest. Weekly growth is especially noisy for rare topics.
Human review is therefore needed before using a rank as evidence of a durable
trend. Source counts should be aggregated and published without user-level data.
The evaluation here measures ranking relevance, not downstream decision quality.
The normalization also discards absolute scale: a small platform with a sharp
relative increase can influence the score as much as a larger platform. This
may be desirable for early detection, but it can amplify noise. In practice,
we would expose both normalized and raw counts, apply a minimum-count
threshold, and let a reviewer inspect source links before acting on a rank.
Those checks would preserve the project's emphasis on traceability.

\section{Conclusion}
The radar demonstrates an inspectable way to combine activity, growth, and
source diversity across platforms. The next step is to test the ranking on
several months of real, legally accessible aggregate data with multiple
independent annotators.

\printbibliography[heading=subbibliography,title={References}]
\end{refsection}
\end{document}

